\chapter{绪论}
绪论（引言）部分应如下几个部分：
\begin{enumerate}
  \item 研究背景：即“这个课题为什么重要？”以及研究领域的历史回顾\notemark{1}；
  \item 文献综述：前人工作的综合述评\notemark{2}；
  \item 研究目的：需要解决什么；
  \item 流程方法：说清楚论文所采用的解决方法流程；
  \item 论文结构：整篇论文的组织架构。
\end{enumerate}

\section{研究背景与意义}

随着智能制造与智慧物流的快速发展，移动机器人已广泛应用于仓储搬运、自动配送、巡检
安防等场景\cite{tan2013robot}。路径规划与避障是其实现自主导航的核心环节，直接决定
机器人能否在狭窄通道、密集人流等复杂条件下安全高效地完成任务。然而，依赖单一传感器
的方案普遍存在局限：激光雷达测距精度高，但在长廊等几何特征退化的环境中容易匹配
失败\cite{zhang2014loam,shan2018lego}；视觉相机成本低、信息丰富，却受光照变化与
纹理缺失影响明显\cite{qin2018vins}；惯性测量单元与轮式里程计输出频率高，但存在
累积漂移\cite{fox2005probabilistic}。因此，本文面向低成本移动平台，研究多传感器
融合的定位方法与分层路径规划算法，以提高机器人在动态环境中的定位精度与避障可靠性。

\section{国内外研究现状}

\subsection{多传感器融合定位研究现状}

早期工作多采用扩展卡尔曼滤波融合惯性测量单元与轮式里程计，随后激光雷达被引入，
形成了以扫描匹配为核心的定位与建图方案\cite{zhang2014loam,shan2018lego}。视觉方面，
多状态约束卡尔曼滤波把视觉观测与惯性状态联合优化\cite{mourikis2007msckf}，
VINS-Mono 进一步实现了单目视觉惯性系统的紧耦合估计\cite{qin2018vins}。国内学者在
机器人感知与导航方面也开展了系统研究\cite{tan2013robot,wang2004parallel}。

\subsection{路径规划与避障研究现状}

路径规划可分为全局规划与局部避障两个层次。全局规划的代表性工作是 A* 算法，它通过
启发式函数引导搜索并保证最优性\cite{hart1968formal}；采样类方法如快速扩展随机树
更适合高维与复杂约束问题\cite{lavalle1998rrt}。局部避障方面，动态窗口法在速度空间
中采样并评价候选轨迹，计算量小、响应快\cite{fox1997dw}。

近年来，深度学习方法被广泛引入感知与决策环节\cite{lecun2015deep,he2016deep}，
深度强化学习则被用于序列决策问题\cite{mnih2015human,silver2016go}，其在未知环境中的
路径规划上表现出一定优势；相关理论基础可参见周志华的专著\cite{zhou2016ml}。

\subsection{存在的问题}

综合国内外研究现状可以看出，现有工作仍存在以下缺口。其一，多数融合方案依赖高精度
传感器与精细标定，在低成本平台上的适用性有待验证；其二，全局规划与局部避障往往
分别设计，缺乏统一的分层框架，难以兼顾路径质量与实时性；其三，动态障碍的运动规律
多被简化处理，与真实场景存在差距。本文即针对上述缺口展开研究。

\section{论文主要研究内容}
本文的主要研究内容包括：
\begin{enumerate}
  \item 构建基于多传感器融合的环境感知模型；
  \item 设计面向动态环境的路径规划与避障算法；
  \item 搭建仿真平台验证所提方法的有效性；
  \item 分析不同场景下的规划效果与鲁棒性。
\end{enumerate}

\section{论文组织结构}
本文共分为六章，各章节内容安排略。